\section{Background}
\label{sec:background-motivation}
\label{sec:motivation-background}

\paragraph{Multi-agent Execution.}
As LLM applications tackle more complex tasks, they increasingly rely on multi-agent collaboration. These applications divide user tasks into subtasks and assign them to several agents for parallel execution. Each Agent maintains its own conversation, advancing tasks through LLM requests and tool invocations: the model generates responses or initiates tool calls, with the returned results becoming input for subsequent requests~\cite{wu2023autogen}. A \emph{request} is a single model call; an \emph{agent} can initiate multiple requests to complete a subtask; and a \emph{workflow} connects these tasks through serial dependencies, parallel branches, and joins. The application is responsible for workflow orchestration, and the inference service layer is responsible for executing model requests.
Figure~\ref{fig:agent-workflow-examples} illustrates this structure. In code generation, the main Agent assigns the CSV export function to the API, UI, test writing, and documentation Agents. The branches execute in parallel, alternating between model requests and file reading, editing, or checking. The main Agent then integrates the changes and runs tests. In deep research, the lead Agent assigns technical, cost, security, and supply investigations, and then combines the collected evidence into a report~\cite{anthropic2025multiagent}. The number of requests and tool calls required by each branch may vary. Therefore, after a model request is completed, its associated Agent and the entire workflow may still not be finished.

\begin{figure}[t]
  \centering
  \includegraphics[width=\columnwidth]{figures/agent-workflow-examples.pdf}
  \caption{Multi-Agent Workflow Examples.
    Each row is one agent: $R_i$ denotes its $i$th LLM request, with
    intervening tool results carried into the next request.
    Branches may have different request counts.
    The coordinator resumes after their results join.}
  \label{fig:agent-workflow-examples}
  \Description{Two stacked left-to-right workflows. A main agent assigns
    CSV-export API, UI, test-authoring and documentation tasks in parallel,
    then integrates their outputs and runs tests. Their request counts
    are four, three, three and two, respectively. A lead researcher assigns
    technical, cost, safety and supply investigations, with four, two,
    three and two requests, before writing a report. Each branch explicitly
    unfolds its requests and intervening tools. Starts and ends are
    staggered; all edges move forward, without loops or round barriers.}
\end{figure}

\paragraph{Workflow completion time.}
Let $w$ represent a workflow, $a_w$ the time the workflow arrives at the system, and $c_w$ the time when all necessary outputs are ready. The total time from arrival to completion of the workflow is defined as the \emph{workflow completion time} (WCT), i.e., $\mathrm{WCT}_w=c_w-a_w$. It includes model execution, tool execution, queuing, and dependency waiting. In the code example, integration needs to wait for all necessary modifications. Delays on the last incomplete branch will delay integration, while branches that are completed earlier may absorb the same delay before merging. Request delays affect the WCT through these dependencies.

\paragraph{KV reuse in multi-agent workflows.}
The model requests on these branches repeatedly consume context from
earlier steps. Coding agents carry source files, edits and test outputs
into later requests; research agents accumulate retrieved documents,
intermediate findings and tool results. Processing this growing history
from scratch on every request would repeat work already performed.
During \emph{prefill}, an LLM computes per-layer keys and values (KV)
for input tokens; \emph{decoding} extends this state as output tokens
are generated. Retaining this KV allows a subsequent request with a
matching token prefix to compute only the new suffix, provided the model
and position semantics are compatible~\cite{gao2024cachedattention,zheng2024sglang}.

This reuse occurs both within and across agents. Successive requests
from one agent can reuse KV for their shared conversation history.
Different agents can also inherit a common prompt before adding their
own tasks. For example, the API, UI and test agents in
Figure~\ref{fig:agent-workflow-examples} can share KV for common project
instructions and repository context when these form the same initial
token sequence. Different role prompts or context edits can shorten
the reusable prefix~\cite{zheng2024sglang}.

The KV cache thus retains intermediate model state accumulated during agent execution and reused across model requests.

\begin{figure}[t]
  \centering
  \includegraphics[width=\columnwidth]{figures/gpu-server-availability.pdf}
  \caption{GPU-server availability reconstructed from InfiniteHBD's
    production-training trace~\cite{shou2025infinitehbd}: 400 servers minus
    those within recorded outage intervals. Records include GPU, network,
    power/cooling, and software faults; unknown causes, tests,
    troubleshooting, and changes are also retained. Overlapping intervals
    per server are merged. Availability is measured at the node level.}
  \label{fig:gpu-server-availability}
  \Description{A step curve over approximately 349 days shows 365 to 400
    GPU servers outside active recorded fault intervals in a fixed
    400-server cohort. The dashed line marks all 400 servers.}
\end{figure}

\paragraph{Failures and KV reconstruction.}
These KV reuse benefits depend on the cached state remaining available.
In large-scale LLM deployments, hardware faults and software crashes can
interrupt serving workers and make their GPU-resident KV
unavailable~\cite{cao2026lumen,jayakody2026ghostserve}.
For production-cluster context, Figure~\ref{fig:gpu-server-availability}
shows the public InfiniteHBD training-cluster trace: 584 recorded outages
affect 231 of 400 GPU servers over approximately 349 days, with reconstructed
availability falling to 365 servers~\cite{shou2025infinitehbd}.
The trace includes faults and operational interruptions, as detailed in
the caption.

When no reusable KV survives, a replacement worker must reconstruct it
by replaying the retained tokens, even if it already has the model weights
and the application has preserved the conversation
history~\cite{strati2024dejavu}. For a long-running agent, this means
reprocessing the context accumulated across earlier requests rather than
only the latest input~\cite{jayakody2026ghostserve}.
One failed worker can also interrupt several co-located agents, requiring
surviving workers to reconstruct multiple histories while handling normal
traffic~\cite{cao2026lumen}. The resulting delays propagate through the
workflow's dependencies: reconstruction on the last unfinished branch
postpones its join and can increase WCT.

\paragraph{KV backup across storage tiers.}
To reduce reconstruction work, a \emph{KV checkpoint} retains an already
computed prefix. If it survives the failure, recovery retrieves this
prefix and replays only the unsaved
tail~\cite{strati2024dejavu,jayakody2026ghostserve}.
KV can be retained in GPU HBM, host DRAM, or shared
storage~\cite{gao2024cachedattention,qin2025mooncake}.
These tiers differ in both storage cost and recovery latency.
HBM offers fast access but has high capacity cost and competes with
foreground inference for GPU memory. DRAM and storage-backed shared
services can provide cheaper capacity, with additional retrieval latency.

Figure~\ref{fig:motivation-recompute} illustrates the latency difference
between reconstructing a history and reusing retained KV.
At 1,010,000 tokens, full prefill takes 571.21\,s even with sparse
attention; HBM, DRAM and shared-storage reuse take 1.305, 2.663 and
10.156\,s, respectively. The history occupies 57.92\,GB across two ranks.
All paths end when the same input's KV is ready: reuse restores the
preceding history, if needed, and computes one new token. These measurements show that reusing retained KV can substantially reduce reconstruction latency, although retaining this state consumes storage capacity.

\begin{figure}[t]
  \centering
  \includegraphics[width=\columnwidth]{figures/recovery-cost.pdf}
  \caption{Measured KV reconstruction and reuse costs for
    Qwen2.5-7B-Instruct-1M on two A100 PCIe 40\,GB GPUs
    (TP2, BF16, batch one), with sparse attention above 32K.
    Full prefill starts from empty KV; reuse starts from $(C-1)$ cached
    tokens and computes one. KV reconstruction time includes any
    retrieval and computation until all $C$ tokens' KV is ready.
    HBM/DRAM whiskers show min--max over three timings; other paths have
    one timing per length. Protocol: Appendix~\ref{sec:motivation-protocol}.}
  \label{fig:motivation-recompute}
  \Description{Four measured KV reconstruction-time curves for Qwen2.5-7B-Instruct-1M
    from 8K to 1,010,000 tokens: blue circles for resident HBM reuse,
    orange squares for DRAM reuse, purple diamonds for shared-storage
    reuse, and red triangles for full prefill. Both axes are logarithmic.
    At the longest context, the four paths take 1.305, 2.663, 10.156 and
    571.21 seconds, respectively.}
\end{figure}

\section{Motivation and Challenges}
\label{sec:motivation}

The resulting management problem is to decide which KV to retain and
in which tier. For a multi-agent service, the benefit of these decisions
depends on how retained state supports both normal execution and workflow
completion after a failure. Existing fault-tolerant serving and
workflow-aware KV management provide useful mechanisms, but their
objectives leave this joint decision unresolved.

\subsection{Limitations of Existing Approaches}
\label{sec:motivation-limitations}

\begin{figure*}[t]
  \centering
  \captionsetup[subfigure]{font=normalsize,labelfont=normalfont,
    textfont=normalfont,justification=centering,singlelinecheck=false}
  \begin{subfigure}[t]{2.48in}
    \vspace{0pt}
    \begin{minipage}[c][1.65in][c]{\linewidth}
      \centering
      \includegraphics[width=\linewidth]{figures/workflow.pdf}
    \end{minipage}
    \caption{Four-branch workflow}
    \label{fig:motivation-workflow}
  \end{subfigure}\hfill
  \begin{subfigure}[t]{1.86in}
    \vspace{0pt}
    \begin{minipage}[c][1.65in][c]{\linewidth}
      \centering
      \includegraphics[width=\linewidth]{figures/workflow-gaps.pdf}
    \end{minipage}
    \caption{Workflow-maximum gaps}
    \label{fig:motivation-branch-gap}
  \end{subfigure}\hfill
  \begin{subfigure}[t]{2.52in}
    \vspace{0pt}
    \begin{minipage}[c][1.65in][c]{\linewidth}
      \centering
      \includegraphics[width=\linewidth]{figures/recovery-exposure.pdf}
    \end{minipage}
    \caption{Recovery cost and WCT}
    \label{fig:motivation-recovery-exposure}
  \end{subfigure}
  \caption{Workflow timing and the WCT impact of KV reconstruction.
    (a) Four trace-derived parallel agents A1--A4 with schematic Main
    agents A0/A5. Bar lengths encode execution time; dashed tails mark
    slack if A5 waits for all four branches.
    (b) Empirical CDF of workflow-level gaps over 121 workflows.
    For each parallel group with temporally inferred fan-in, we subtract
    the shortest agent execution time from the longest, then retain the
    largest gap per workflow. These gaps are not failure-time slack.
    (c) Reconstruction costs for A1--A4 in (a), measured by full prefill
    at trace-estimated context lengths using
    Figure~\ref{fig:motivation-recompute}'s setup. Total bar lengths and
    whiskers show the median and min--max over three timings after warmup.
    Each measured cost is added as a pause to its corresponding branch,
    one branch at a time, with other execution fixed. If A5 waits for all
    branches, the hatched portion is absorbed by slack and the solid
    portion is the calculated WCT increase.
    Protocol: Appendix~\ref{sec:motivation-protocol}.}

  \label{fig:motivation}
  \Description{Left: a four-branch workflow with A2 completing last at
    151.6 seconds. Middle: an ECDF of workflow-maximum branch-duration
    differences with a median of 138 seconds and a 90th percentile of
    1173 seconds. Right: A1 through A4 have measured median Prefill costs
    of 33.2, 32.6, 35.4 and 39.9 seconds. Applying each cost separately
    as an added branch pause gives calculated WCT increases of 0.0, 32.6,
    16.3 and 18.0 seconds. Hatching identifies the absorbed portion;
    whiskers show timing minima and maxima, not confidence intervals.}
\end{figure*}

\paragraph{Limitation 1: Fault tolerance without workflow-aware priorities.}
Fault-tolerant serving systems preserve state and coordinate recovery
for interrupted requests. D\'ej\`aVu streams KV checkpoints, and GhostServe
uses host-memory parity to reduce reconstruction
work~\cite{strati2024dejavu,jayakody2026ghostserve}.
LUMEN further coordinates checkpoint placement, recovery routing and
capacity restoration according to worker load~\cite{cao2026lumen}.
These mechanisms address state loss and recovery contention, but do not
use multi-agent dependencies to determine which requests' faster
recovery would advance workflow completion. Request recovery cost and
worker load alone cannot express this benefit.

\label{sec:motivation-workflow}
Figure~\ref{fig:motivation}(\subref{fig:motivation-workflow}) illustrates
why workflow position matters. A2 finishes last, 151.6\,s after the fork,
whereas A1 finishes at 60.4\,s. If the following Main stage requires all
four outputs, A1 can absorb 91.2\,s of additional delay before postponing
the join, whereas any delay on A2 postpones it. We call the additional
delay a branch can absorb without postponing the required output its
\emph{slack}.

Such timing differences occur across the retained CC/Weka
traces~\cite{semianalysis2026cctraces}. From 393 filtered Claude Code
sessions, a temporal fan-in proxy retains 286 overlapping subagent
groups in 121 workflows: all branches complete, and a Main request
follows the last completion; Main activity during the branches is allowed.
Figure~\ref{fig:motivation}(\subref{fig:motivation-branch-gap}) reports
each workflow's maximum longest-minus-shortest branch-duration gap,
with a median of 138\,s and a 90th percentile of 1,173\,s.
The trace lacks explicit result-consumption edges, so these differences
neither verify all-branch joins nor measure failure-time slack.
The four-branch example is one retained group, not a representative
workflow-maximum gap.

To connect reconstruction cost to workflow delay, we measure full
prefill at the four active requests' trace input-size proxies, using
Figure~\ref{fig:motivation-recompute}'s Qwen2.5-7B-Instruct-1M
configuration with sparse attention~\cite{qwen2025qwen25oneM,jiang2024minference}.
Figure~\ref{fig:motivation}(\subref{fig:motivation-recovery-exposure})
then inserts each measured cost $R_i$ separately as a pause on branch
$i$ at 29.5\,s after the fork. With offline slack $\sigma_i$, an assumed
all-branch join, and all other execution and queues fixed, the WCT
increment is
\[
  \Delta\mathrm{WCT}_i=\max(0,R_i-\sigma_i).
\]
A1's 33.2\,s reconstruction fits within its 91.2\,s slack and leaves WCT
unchanged. A2's slightly smaller cost, 32.6\,s, is fully exposed because
its slack is zero. Thus prioritizing the more expensive reconstruction
can favor a branch with less immediate effect on completion.
Only the Prefill costs are measured; WCT increments are conditional
calculations using synthetic inputs and trace timing, not fault-injected
workflow measurements. Appendix~\ref{sec:motivation-protocol} gives the
input construction, timing samples and mapping protocol.

\paragraph{Limitation 2: Workflow-aware caching without a recovery objective.}
Workflow-aware KV management uses application structure to improve
normal execution. KVFlow uses upcoming agent invocations to guide prefix
retention and prefetching~\cite{pan2025kvflow}; TOPAS jointly schedules
prefix residency and request admission using remaining workflow paths,
reuse benefits and movement or preemption costs~\cite{ni2026topas}.
These systems already exploit workflow information, but their cache
decisions do not explicitly optimize protection against serving-worker
failures: which state must survive, where to retain it, and how quickly
it should become available to a replacement worker.

The distinction changes the value of retained state. For example,
KVFlow gives varying suffixes the highest eviction priority while
protecting reusable fixed prompts~\cite{pan2025kvflow}. A suffix with
little expected reuse can nevertheless be needed to resume an
interrupted execution. Likewise, keeping a prefix in the executing
worker's HBM accelerates normal access but does not by itself preserve
that state when the worker fails. A cache copy can support recovery if
it survives; its usefulness for future calls alone does not determine
its value as protection against a particular failure.

\label{sec:motivation-long-context}
The tier costs in Figure~\ref{fig:motivation-recompute} make the resulting
allocation tradeoff concrete. Retaining KV in a slower tier already
avoids much of the full-prefill cost; faster tiers reduce the remaining
retrieval delay at a higher capacity cost. Whether that reduction
advances workflow completion depends on the recovering branch's slack.
Conversely, a branch with ample slack may still reuse its history during
normal execution, so low recovery urgency does not justify deleting its
last reusable copy. These are distinct reasons to retain and accelerate
state, even when both use the same storage hierarchy.

\subsection{Design Challenges}
\label{sec:motivation-challenges}

Together, these limitations call for a policy that connects workflow
progress, failure recovery and KV storage management. Closing the gap
requires estimating the value of retained state, translating that value
into placement decisions, and maintaining those decisions as execution
changes. These are coupled challenges: the usefulness of fast recovery
depends on workflow progress, while protection itself consumes resources
needed by the workflow.

\paragraph{Challenge 1: Estimating the value of KV protection online.}
A protection decision must relate saved reconstruction or retrieval time
to the workflow's required outputs. The example above obtains slack from
completed traces; an online system must act before the remaining model
and tool durations are known. The branch that gates completion can
change, and one worker failure can interrupt several agents whose
recoveries share resources. Normal reuse also contributes value even
when recovery is not urgent. The challenge is to derive useful priorities
from partial execution and reuse history without treating request-level
savings as independent, additive WCT reductions.

\paragraph{Challenge 2: Translating protection value into tiered placement.}
A request's urgency alone does not specify how much of its KV should
occupy each tier. A policy must distinguish retaining enough state to
avoid reconstruction from using faster capacity to reduce retrieval
latency. It must compare these benefits with the bytes retained, the
capacity cost of each tier and the work needed to move state.
The same KV may support several agents, while different portions of a
long history may have different access demands. Allocating capacity
therefore requires connecting workflow-level priorities to the state
that recovery and reuse actually need, subject to finite tier budgets.

\paragraph{Challenge 3: Adapting protection without disrupting inference.}
Agent progress, tool returns and new requests change both protection
value and foreground memory demand. Repeatedly reassessing every stored
KV object can make fine-grained management expensive; frequent placement
changes can also consume bandwidth needed by active requests.
A useful policy must update protection with bounded decision and
migration work, react to capacity pressure, and retain valid recovery
paths while state moves between tiers. The challenge is to keep
protection responsive without allowing its management and data movement
to erase the workflow benefit it seeks.

\label{sec:motivation-goals}
These challenges motivate AegisKV, a workflow-aware KV protection system
that coordinates recovery priorities, tiered placement and online state
management. The next section develops the observations and mechanisms
behind this design.
